\runninghead{Friend, How Camest Thou In?}
\properlane{wedding-garment}{Friend, How Camest Thou In? The Garment the Called Must Wear}
\label{sec:the-wedding-garment}

Where the longer form of the Gospel is read, the parable does not end with the
hall full. The king comes in to see the guests, finds one without a wedding
garment, and asks him, ``Friend, how camest thou in hither not having on a
wedding garment?'' The man is silent, and is bound and cast out, ``For many
are called, but few are chosen'' (Mt 22:11--14). The difficulty is plain. The
man did not force his way in: the king's own servants gathered him from the
roads with everyone else, ``both bad and good'' (22:10). Why is he blamed for
his clothes? This interpretation answers with what the Fathers who expound
these verses share, each in his own words: the call is free and goes out to
all, and the garment is required of those whom the call has gathered, is to be
kept, and is answered for when the king comes in. Those who say where the
garment comes from, Chrysostom, Augustine, Gregory and Hilary, say that it is
given, as the call is. Irenaeus and Hilary state the freedom of the call;
Chrysostom and Aquinas say why the called are nonetheless examined. What the
garment is, they say differently. Augustine and Gregory name it charity, and
Aquinas counts charity among the ways of putting on Christ; Chrysostom calls
it life and practice, Irenaeus works of righteousness, Jerome the Lord's
precepts and the works of the new man, and Hilary the glory of the Holy
Spirit. The interpretation leaves that identification open and sets the
answers side by side. The
joins among them, and to the Collect and the rest of the Mass, are the
editor's. The Gospel at the centre of this interpretation stands only in the
longer form; what remains of it under the shorter form is said at the end.

\subsection{Called without exception, and examined}

The witnesses first state the freedom of the call. Irenaeus says that the one
God ``does indeed, through His infinite kindness, call the unworthy; but He
examines those who are called'', whether they have on the garment fit for his
Son's marriage (\work{Against Heresies} IV.36.6). Hilary of Poitiers puts the
call's holiness and the guests' failure in one sentence: \latin{Vocatio quidem
bonos efficere debuerat, quia sancta est, et ex optimo affectu invitantis
profecta est: sed per vitium inemendatae voluntatis discrimen est vocatorum},
the calling ought to have made them good, because it is holy and came from the
best affection of the one inviting; but through the fault of an unamended will
there is a difference among the called (\work{Commentarius in Matthaeum} 22.6,
PL 9, col.~1043). He also asks the reader's question himself. Had the king
laid down what the invited were to wear? He had ordered whoever was found to
be invited, so how could all have been dressed alike? And if a fixed dress had
been the custom at weddings, the servants could have kept the man out. His
answer is that the fault was one only God could see: \latin{idcirco hunc malum et
indignum coetu nuptiali Deus solus invenit}, therefore God alone finds this man,
evil and unworthy of the wedding company (22.7, col.~1044).

John Chrysostom gives the clearest account of how the call and the garment
belong together. Having called the gentiles, the Lord ``discourses unto them
also concerning the judgment to be passed upon wicked actions'', so that no one
should ``put confidence in their faith alone''. ``For the garment is life and
practice. And yet the calling was of grace; wherefore then doth He take a
strict account? Because although to be called and to be cleansed was of grace,
yet, when called and clothed in clean garments, to continue keeping them so,
this is of the diligence of them that are called'' (\work{Homilies on Matthew}
69.2). In Chrysostom's account, then, the guest received a clean garment with
his call and did not keep it. Thomas Aquinas meets the objection directly:
someone might say the king had occasion against the man, since he had called
the bad as well as the good; but \latin{noluit quod mali venirent nisi pararent
se, et disponerent se ut essent boni}, he did not will that the bad should come
unless they prepared themselves and disposed themselves to be good
(\work{Super Matthaeum} c.~22, Venice, p.~282).

\subsection{Charity: the answer of Augustine and Gregory}

Augustine searches for the garment by elimination, twice. In one sermon he
asks whether it is a sacrament: ``Is it Baptism? Without Baptism it is true no
one attaineth to God; but not every one that hath Baptism attaineth to Him. I
cannot therefore understand Baptism, the Sacrament itself that is, to be `the
wedding garment;' for this garment I see in the good, I see in the bad.'' Nor is
it the altar, for many ``eat and drink judgment to themselves''; nor fasting,
nor running to church, nor miracles. ``What is that `wedding garment' then?
This is the wedding garment: `Now the end of the commandment,' says the
Apostle, `is charity out of a pure heart, and of a good conscience, and of faith
unfeigned'\,'' (1~Tim 1:5). Charity in the guest is something living: ``Let
charity be born in thee, if it be yet unborn, and if it be born, be it
nourished, fostered, increased'' (\work{Sermo} 90.5--6). In the other sermon he
lists again what the good and the bad share, baptism, ``the Sacraments of the
Altar'', prophecy and faith, for ``The devils also believe and tremble'', and
finds the garment only in Paul's ``If I have not charity''. He takes up the
objection of the excessive punishment: ``For so small a fault, so great a
punishment? For great it is. It is called a small fault not to have `the
wedding garment;' small, but only by those who do not understand.'' And he says
where the garment comes from: ``Do not say; `we are too poor to have that
garment.' Clothe others, and ye are clothed yourselves. It is winter, clothe
the naked. Christ is naked; and He will give you that `wedding garment'
whosoever have it not. \ldots\ He knoweth how to clothe His naked ones''
(\work{Sermo} 95.5, 7).

Gregory the Great reaches the same answer by the same route. If the garment
were baptism or faith, he asks, who could have come into this wedding without
them? \latin{Quid ergo debemus intelligere nuptialem vestem, nisi
charitatem?} What then are we to understand by the wedding garment, except
charity? It is rightly called the wedding garment, he says, because the
Creator himself had it \latin{dum ad sociandae sibi Ecclesiae nuptias venit},
when he came to the wedding in which he joined the Church to himself; it was by
love alone that the Only-begotten united the minds of the elect to himself. The
believer has entered the wedding, but has not come in the wedding garment,
\latin{si charitatis gratiam non custodit}, if he does not keep the grace of
charity (\work{Homiliae in Evangelia} 38.9). The garment is woven, like any
cloth, between two beams, and charity likewise holds in two commandments, the
love of God and the love of neighbour (38.10); and the love of neighbour is
proved by the love of enemies (38.11).

For Augustine and Gregory, then, the garment is not something the bad can
have, and that is why it cannot be baptism or faith. Both say that it is given
and must be kept: Augustine's Christ ``will give you that `wedding garment'\,'',
and Gregory's guest lacks it if he does not keep \latin{charitatis gratiam}.
On their answer the difficulty of the guest from the roads is met by what the
garment is: the call gathered him, and charity was what he had not put on.

\subsection{Other names for the garment, and a real disagreement}

The other witnesses name the garment in their own terms. For Chrysostom it is
``life and practice'', the clean garment received with the call and kept by
diligence. For Irenaeus, ``we ought, after our calling, to be also adorned with
works of righteousness, so that the Spirit of God may rest upon us; for this is
the wedding garment''; he cites Paul's wish to be ``clothed upon, that
mortality might be swallowed up by immortality'' (2~Cor 5:4). For Jerome it is
the new man's clothing: \latin{Vestis autem nuptialis, praecepta sunt Domini,
et opera quae complentur ex lege et Evangelio, novique hominis efficiunt
vestimentum}, the wedding garment is the Lord's commandments and the works that
are done from the law and the Gospel and make the new man's clothing; the man
cast out wears instead \latin{veteris hominis exuvias}, the cast-off skin of the
old man (\work{Commentarii in Matthaeum} III, PL 26, cols.~160--161). Aquinas
gathers these into one: \latin{Quae est ista vestis? Christus. Qui sumus
Christi, Christum induamus}, what is this garment? Christ: we who are Christ's,
let us put on Christ (Rom 13:14). Some put him on by the sacrament (Gal 3:27),
some by charity and love (Col 3:14), some by remembering death, some by
conforming their works to his; \latin{Habere ergo vestem nuptialem est induere
Christum per operationem bonam, per conversationem sanctam, per caritatem
veram: et si unum deficiat, malum}, to have the wedding garment, then, is to
put on Christ by good works, by holy living, by true charity, and if one is
lacking, it is bad (p.~282).

Hilary of Poitiers gives an answer that stands against Augustine's and
Gregory's. \latin{Vestitus autem nuptialis est gloria Spiritus sancti, et candor
habitus coelestis: qui bonae interrogationis confessione susceptus, usque in
coetum regni coelorum immaculatus et integer reservetur}: the wedding garment is
the glory of the Holy Spirit and the whiteness of the heavenly habit, which,
received at the confession of the good questioning, is to be kept unstained and
whole until the assembly of the kingdom of heaven (22.7, col.~1044). Hilary
does not say baptism. But the phrase ``the good questioning'' recalls Peter's
words on baptism, which the Clementine Vulgate gives as \latin{conscientiae
bonae interrogatio in Deum}, ``the examination of a good conscience towards
God'' (1~Pet 3:21); read as an allusion to that verse, which is the editor's
reading and not Hilary's statement, his garment is received at baptism and kept
by the baptized. The \work{Procatechesis} of Cyril of Jerusalem uses the parable
in the same setting: the man who came in ``when he saw them all clad in white,
\ldots\ ought to have assumed a garment of the same kind himself'', and Cyril
tells his candidates for baptism, ``Thou seest what happened to that man: make
thine own condition safe'' (3).

This disagreement is real and it matters. Augustine and Gregory say expressly
that the garment is not baptism, because the bad are baptized; Hilary's
garment, on the reading of his allusion given here, is the gift received at the
baptismal confession and kept. Chrysostom's clean garment given with the call
and kept by diligence stands between them, and Aquinas includes the sacrament
among the ways of putting on Christ while insisting that charity and works
must be there too. The witnesses agree on more than they dispute: all of them
make the garment something required of those whom the call has already
gathered and something to be kept, and all who comment on the king's coming
in place it at the judgment. That is the claim this interpretation makes. On
what the garment names they do not agree, and the interpretation does not
decide between them. Charity, works, Christ put on and the Spirit's glory are
each given as their authors' answers, and Augustine's and Gregory's denial
that the garment is baptism stands against Hilary on the reading of his phrase
given above.

\subsection{``Friend'', and the silence}

The king's question begins with a word that the witnesses weigh. Jerome:
\latin{Amicum vocat, quod invitatus ad nuptias est}, he calls him friend because
he was invited to the wedding (col.~161). Gregory: \latin{Amice, et non amice;
amice per fidem, sed non amice per operationem}, friend and not friend, friend
by faith but not friend in deed (38.12). The man's silence is the end of all
excuse: \latin{in illa districtione ultimae increpationis omne argumentum
cessat excusationis}, in the strictness of that last rebuke every argument of
excuse falls silent, because the one who rebukes from outside is the witness
of the conscience within (38.12). Aquinas: \latin{non potest habere
sufficientem rationem peccator quare vestem nuptialem contempsit}, the sinner
can have no sufficient reason why he despised the wedding garment (p.~282).

Gregory at once tempers the scene. Whoever has this garment of virtue, but not
yet perfectly, \latin{ad pii regis ingressum desperare veniam non debet}, must
not despair of pardon when the merciful king comes in, since the Psalmist says,
``Thy eyes did see my imperfect being, and in thy book all shall be written''
(Ps 138:16; 38.12). The king's coming in is the judgment. Jerome says he comes
\latin{ut in die iudicii visitare convivas, et discerneret merita singulorum},
in order on the day of judgment to visit the guests and discern the merits of
each (col.~160). Aquinas says the king comes in when he
exercises judgment, at the final judgment, and also at death and when
tribulations press on the Church (p.~282). Gregory says that the king ``is to
enter'' on the coming day of judgment, \latin{futuro die iudicii rex
intraturus est} (38.14).

\subsection{Many are called}

The witnesses who comment on the closing saying weigh how few the chosen are,
and they reach that judgement from different places. Gregory passes to it from
the guest cast out: \latin{repulso uno, in quo videlicet omne malorum corpus
exprimitur}, the one rejected, in whom the whole body of the wicked is
expressed, is followed at once by the general sentence (38.14). Augustine makes
the same move and states the proportion: ``Leave alone the few, cast out the
many. It is true, that man was but one. Yet undoubtedly that one not only was
many, but those many in numbers far surpassed the number of the good. For the
good are many also; but in comparison of the bad, they are few.'' He says it
twice, ``The same persons considered in themselves are many, in comparison with
the bad are few'', and says what the saying shows, ``who in this present feast
are accounted to be such, as to be brought to that other feast, where no bad
men shall come'' (\work{Sermo} 90.4). At the feast now spread, then, the bad
far outnumber the good, and the good, many in themselves, are few beside them.

Gregory had already said as much of the hall of bad and good, which for him is
the present Church. His hearers are not to be alarmed \latin{quod in Ecclesia
et multi mali et pauci sunt boni}, that in the Church the bad are many and the
good few. The ark of the flood, a figure of the Church, was broad below, where
it carried the beasts, and narrowed above to a single cubit, where it kept the
men, because the holier men are, the fewer; ``wide is the gate, and broad is
the way that leadeth to destruction, and many there are who go in thereat'',
and few find the narrow way that leads to life (Mt 7:13--14); and on the
threshing floor \latin{pauca sunt grana quae servantur horreis, et grandes
acervi palearum qui ignibus comburuntur}, few are the grains kept for the barns,
and great the heaps of chaff burned in the fire (38.8). Aquinas gives the saying
its reason in the parable and joins it to the same verse of Matthew 7:
\latin{quia quidam nolunt venire, quidam non habent vestem nuptialem}, because
some will not come and some do not have the wedding garment; whence, he says,
Matthew 7:14, ``strait is the way that leadeth to life'', \latin{et pauci sunt
qui inveniunt eam}, ``and few there are that find it'' (p.~282). Irenaeus reads it beside the
former covenant: ``For as in the former covenant, `with many of them was He not
well pleased;' so also is it the case here, that `many are called, but few
chosen'\,'' (IV.36.6; 1~Cor 10:5).

Hilary puts the scarcity among the chosen and not among the invited:
\latin{Non est igitur paucitas in invitatis, sed raritas in electis: quia in
invitante sine exceptione publicae bonitatis humanitas est; in invitatis vero
de iudicii merito probitatis electio est}, there is no fewness among the
invited, but a rarity among the chosen, because in the one inviting there is a
kindness of public goodness without exception, while among the invited the
choosing is of their uprightness, by the merit of the judgment (22.7). The
same witnesses keep the call itself without limit: Hilary in that sentence;
Augustine in saying that the Lord's table is spread for all who will (90.1);
Irenaeus in the King who ``gathered from all quarters the faithful to the
marriage of His Son'' (IV.36.6); Gregory in the Church that now receives bad
and good without distinction and parts them only at the going out (38.7).
Jerome reads the saying differently, as the sentence that sums up every parable
of this kind, the vineyard, the house and the wedding: \latin{non initia, sed
finis quaeratur}, not the beginning but the end is sought (col.~161).

Gregory turns the saying on his hearers. \latin{Quia vocati sumus, novimus; si
sumus electi, nescimus}, that we are called we know; whether we are chosen we
do not know. \latin{Tanto ergo necesse est ut unusquisque nostrum in humilitate
se deprimat, quanto si sit electus ignorat}: so each of us must humble himself
as much as he does not know whether he is chosen (38.14). He ends the homily on
the same note: since no one is certain of his election, \latin{omnes in sola
divina misericordia gaudeant, nullus de suis viribus praesumat}, let all
rejoice in the divine mercy alone, and let no one presume on his own strength
(38.16). The fewness he has affirmed of the Church becomes, for each hearer,
humility about himself and hope in mercy.

\subsection{Grace that goes before and follows}

The Collect asks for exactly what this interpretation says the guest lacked:
that grace go before and follow the assembly and keep it intent on good works.
No witness joins the prayer to the parable; but the pair of verbs it uses has a
history in the reading of this Sunday's psalm. Aquinas, commenting on the
psalm's last verse, ``thy mercy will follow me'', writes: \latin{Alias petivit ut
praeveniat, hic quod subsequatur: et utraque est necessaria: quia praeveniens
est necessaria, quia inspirat animum, subsequensque iuvat ut efficiatur}:
elsewhere the Psalmist asked that it go before, here that it follow, and both
are necessary, the prevenient because it inspires the mind, and the subsequent
because it helps the work to be done (\work{Postilla super Psalmos} 22). In the
\work{Summa} he founds the division of grace on the same verse and on ``his mercy
shall prevent me'' (Ps 58:11), and quotes Augustine: \latin{praevenit ut
vocemur, subsequitur ut glorificemur} (I-II, q.~111, a.~3). The psalm verse
shares one of the Collect's two verbs, \latin{subsequetur} beside
\latin{sequatur}, and the pair of going before and following is Aquinas's, who
joins Psalm 58:11 to this verse. Heard with the parable, this is the grace
that clothes the guest and keeps the garment on him.

The commentators on the prayer, quoted with the Collect above, say what that
grace does: without its going before, Fromage writes, there is not even the
thought of good, and without its following no act of virtue; ``If, on the other
hand, we be faithful to grace, our life will be one uninterrupted tissue of
good works''. Neither he nor Schuster joins the prayer to the parable.
Schuster, commenting on the Gospel itself, gives the parable's lesson in a
sentence that holds the call and the garment together: ``God calls our souls,
but they must be in accord with their calling, so that the grace of eternal
happiness may also be their due reward'' (\work{The Sacramentary}, vol.~3,
p.~173).

\subsection{The Missal's texts in this reading}

The Entrance Antiphon asks who could stand if the Lord marked iniquities, and
in this interpretation it hears the guest's silence before the question was
asked of him. Bellarmine's reading of the verse, given with the antiphon
above, that without grace we can neither make satisfaction nor even see the
enormity of the offence, is the reason no one could answer the king by his own
resources, and Augustine's propitiation, the sacrifice offered for us, is the
hope on which the antiphon turns. The first reading enters by a
word. Isaiah's death cast down for ever is quoted by Paul as ``Death is
swallowed up in victory'' when the mortal has put on immortality (1~Cor
15:53--54), and Irenaeus, explaining the garment, cites Paul's wish to be
``clothed upon, that mortality might be swallowed up by immortality''. The
three passages share the swallowing of death, and in this interpretation the
garment is what the mortal puts on when death is swallowed up.

The second reading names the strength in which the garment is worn. On ``I can
do all things in him who strengtheneth me'' Chrysostom and Aquinas, quoted
under the third interpretation, both give the strength to God and not to Paul;
the third \work{Mystagogical Catechesis} puts the verse in the mouth of the
newly anointed. The acclamation prays for the enlightened eyes that know the hope of
the calling, and Chrysostom's comment that God ``would thus have us become good
of our own will'' is this interpretation's whole difficulty in a sentence: the
call is God's and the response is willed. The Prayer over the Offerings asks
that through these acts of devout worship, \latin{per haec piae devotionis
officia}, the faithful pass to heavenly glory; heard here, the works of which
the garment is made are offered.

Where the first Communion antiphon is chosen, the rich who went hungry are set
against the heart that Augustine calls ``filled with such jewelry of virtues,
justice, truth, charity, faith, endurance'' (\work{Enarrationes in Psalmos}
33). Where the second is chosen, the promise of seeing him as he is has, in
the letter, the next verse beside it: ``every one that hath this hope in him
sanctifieth himself'' (1~Jn 3:3), which the antiphon does not include.
Augustine's comment on that verse joins grace and will: ``Who purifieth us but
God? Yea, but God doth not purify thee if thou be unwilling. \ldots\ Thou
purifiest thyself, not by thyself, but by Him who cometh to inhabit thee''
(\work{Homilies on the First Epistle of John} 4.7). And Aquinas makes charity
the measure of the vision itself: ``he will have a fuller participation of the
light of glory who has more charity'' (\work{Summa theologiae} I, q.~12, a.~6).
The Prayer after Communion asks to be made partakers of the divine nature.
Aquinas calls grace that participation and derives the infused virtues from
it (I-II, q.~110, a.~3); in this interpretation the garment is the life of that
grace. The verse of 2~Peter whose words the
prayer uses goes on, ``flying the corruption of that concupiscence which is in
the world''.

\subsection{Where this interpretation strains}

The first limit is the Gospel's form. Under the shorter form the reading ends
at v.~10, and of this interpretation's Gospel centre only the gathering of
``bad and good'' remains, which Augustine reads as the present feast of good
and bad and Gregory as the present Church where they mingle. The free call
survives; the garment does not.

The second is that every answer to the guest's difficulty is its author's and
not the parable's. The parable does not say what the garment is, or where it
was to be had. Augustine answers that the fault is not small because the
garment is charity; Aquinas, that the king did not will the bad to come unless
they prepared themselves; Hilary, that God alone saw what the servants could
not; Cyril, that the guest saw the others in white.

The third is the disagreement already set out: the witnesses do not agree on
what the garment is, and the interpretation, which rests on what they share,
cannot settle it; whoever names the garment charity speaks with Augustine and
Gregory, and not with Hilary.

Finally, the closing saying is larger than the scene it closes. The parable
shows one guest cast out of a full hall; the saying speaks of many called and
few chosen. Augustine and Gregory make the one stand for many, Gregory for the
whole body of the wicked; Gregory and Aquinas read the few with the strait way
that few find, Gregory adding the ark and the threshing floor, and Irenaeus
with the many in whom God was not well pleased under the former covenant;
Hilary finds the rarity among the chosen and never in the call. Their judgement
on how few are chosen is drawn from those texts. The saying itself names no
number.

\Needspace{26\baselineskip}
\subsection{The four senses of this interpretation}

\begin{fourSenses}
\item[Literal.] In the longer form of the Gospel, the king, coming in to see
the guests his servants gathered, finds one without a wedding garment, who has
nothing to say and is cast out, for many are called and few are chosen. Under
the shorter form, the servants gather bad and good alike, and the hall is
filled.

\item[Allegorical.] The garment is charity, which the Creator himself had when
he came to the wedding of the Church (Gregory); or Christ put on by sacrament,
charity and works (Aquinas); or the glory of the Spirit received at the
confession and kept to the kingdom (Hilary). Grace goes before, calling, and
follows, bringing to glory (Aquinas, with Augustine as he quotes him).

\item[Moral.] Let charity be born, nourished and increased; clothe the naked,
and be clothed (Augustine). Keep the garment the call gave you (Chrysostom).
Be humble, not knowing whether you are chosen, and do not despair of a garment
still imperfect (Gregory).

\item[Anagogical.] In the longer form, the king's coming in is the judgment
(Jerome, Gregory), at
which not the beginning but the end is sought (Jerome), and few are chosen,
as few find the strait way that leads to life (Gregory, Aquinas). Where the second
Communion antiphon is chosen: the more charity, the fuller the vision of him as
he is (Aquinas).
\end{fourSenses}
